\section{Related work}
\label{sec:related}

We situate \pqlang{} among quantum programming languages and frameworks,
then among fault-tolerant algorithm-analysis libraries, and finally among
the commercial platforms that now target quantum scientific computing; the
feature comparison of Table~\ref{tab:qpl-comparison}, the coverage
comparison of Table~\ref{tab:ode-coverage}, and the platform comparison
of Table~\ref{tab:platform-comparison} record the outcome.  The
discussion is organized by the
requirements R1--R8 of \secref{sec:qsc-requirements}; general surveys of the
field are available~\cite{selinger2004towards,heim2020quantum,garhwal2021systematic}.

\subsection{Quantum programming languages}
\label{sec:related-qpl}

\paragraph{Functional and typed language lines.}
Quipper~\cite{green2013quipper} pioneered scalable, circuit-generating
functional programming: a Haskell host generates gate-level circuits, with
``boxing'' providing a form of abstraction over generated subcircuits.
Scaffold~\cite{javadiabhari2015scaffcc} embedded modules and classical
control in a C-like surface, compiled by ScaffCC with early
resource-estimation passes.  LIQUi|>~\cite{wecker2014liquid} was the
earliest vertically integrated attempt: an F\#-embedded DSL with
compilation, simulation, and gate-count resource costing, whose authors'
repository later shipped a complete phase-estimation-based linear-system
example (QLSA)---both our workload and our resource axis appeared there,
but as monolithic, concrete circuits with no oracle abstraction; the
project is unmaintained and was superseded by Q\#.  Notably, Quipper later served as the vehicle
for landmark resource analyses of quantum linear systems---although as
external studies over generated circuits, not as a language
mechanism~\cite{scherer2017resource,fillion2017dirac}.  QWIRE and the Proto-Quipper line developed linear and dependent type
disciplines for safety guarantees~\cite{paykin2017qwire,fu2023protoquipper}.
Silq~\cite{bichsel2020silq} introduced safe automatic uncomputation through
its type system---a genuine language-level answer to the invertibility half
of our R4.  On the standalone-language axis, isQ~\cite{guo2023isq}
integrates a compiler and instruction set with classical control in a
maintained software stack; its public compiler repository has been
inactive since 2023, a data point on the sustainability cost of
standalone language efforts.  None of these lines, however, provides deferred-oracle
semantics with paradigm contracts and late linking (R1/R2): abstractions
expand to circuits, and the open state of a program is not a shareable
artifact.  Their IRs target gate-level compilation rather than
module-preserving program storage (R3/R5), and none ships a
scientific-computing algorithm stack.

\paragraph{Industrial SDKs and embedded frameworks.}
ProjectQ~\cite{steiger2018projectq} established the Python-embedded pattern,
shipped one of the earliest resource estimators, and---through its
extensible compiler-engine chain---explicitly supports \emph{emulating}
quantum programs at a high level of abstraction, mimicking large oracles
instead of compiling them to gates: a genuine form of deferred oracle at
execution time, though one that lives in the engine stack rather than in a
serializable open program.  Quil and pyQuil~\cite{smith2016quil} and
Cirq~\cite{cirqdevelopers2020cirq} provide circuit DSLs centered on NISQ execution.
Qiskit~\cite{javadiabhari2024qiskit} grew a large circuit library and textbook algorithm
implementations; demonstration-level HHL shipped historically in its
algorithms module and has since been retired with it, and no QSVT or
block-encoding abstraction exists in the library today.  Its objects are
concrete circuits, its transpiler flattens structure, and alternative oracle
realizations are separate programs rather than bindings of one open slot.
PennyLane~\cite{bergholm2018pennylane} centers differentiable hybrid
programming and, notably, ships subroutine templates that overlap our
primitive layer---\texttt{QSVT}, \texttt{GQSP}, \texttt{Qubitization}, and
the \texttt{FABLE} block-encoding construction---together with an
experimental fault-tolerant computing module; these are circuit templates
with externally supplied phase angles, not packaged solvers, and no QLSS
or ODE/PDE stack is assembled from them.  The mainstream Chinese
SDKs---QPanda/pyQPanda~\cite{originq2024qpanda} and MindSpore
Quantum/MindQuantum~\cite{xu2024mindquantum}---are NISQ- and
variational-centric, with at most demonstration-level HHL examples and no
block-encoding or QSVT abstractions.
OpenQASM~3~\cite{cross2022openqasm3} standardizes a textual interchange
language whose subroutines and loops preserve some structure, and whose
``opaque'' declarations gesture toward R1---but at gate granularity,
without paradigm contracts, capabilities, or a linking discipline.
Closest to our primitive layer among embedded frameworks, Qrisp~\cite{seidel2023qrisp}
automates uncomputation and resource-friendly scheduling inside Python and
exposes a \texttt{BlockEncoding} class---block-encodings as programming
abstractions with a {NumPy}-like interface---which we regard as independent
evidence that scientific-computing primitives belong at the language level.
Qrisp does not, however, treat the open program as an artifact: there is no
paradigm-contracted oracle declaration with late, partial, repeatable
linking, no module-preserving serialized IR with promise provenance, and no
ODE/PDE solver families.  t\textbar ket\textgreater~\cite{sivarajah2021tket} and
staq~\cite{amy2020staq} are compiler and synthesis toolchains,
complementary to our concerns: they operate \emph{below} the register/module
level our RIR preserves.

\paragraph{Q\#.}
Q\#~\cite{svore2018qsharp} is the closest prior art at the language level.
Its operations with \texttt{Adjoint}/\texttt{Controlled} \emph{functors}
anticipate our capability discipline (R4): specialization availability is
declared and tracked.  Passing operations as parameters---the closest
mainstream analogue to our oracle slots---lets algorithms be written against
any conforming implementation, so Q\# partially meets R1 at \emph{call}
time.  Its standard libraries include fixed-point numerics and estimation
primitives, and the surrounding Azure tooling ships a serious resource
estimator.  What Q\# does not provide for our workload is
the \emph{open program}: oracles are operations with concrete bodies;
``partial application'' is host-level; there is no serialized artifact in
which an algorithm carries its unimplemented input models with paradigm,
$\alpha$, capabilities, and provenance (R1/R2/R6); and its libraries do not
include block-encoding algebra, a QSVT/QSP engine, or any ODE/PDE solver
families (Table~\ref{tab:ode-coverage}).  The distinction is one of
\emph{when} semantics attach: Q\# resolves capabilities at compile time of
concrete operations, whereas \pqlang{} keeps the \emph{link} itself as a
first-class, repeatable, partial operation on stored programs.

\paragraph{Fault-tolerant algorithm-analysis libraries.}
In the open-source world, the systems closest to our primitive layer are
libraries for compositional implementation, simulation, and resource
analysis of fault-tolerant algorithms.
Qualtran~\cite{harrigan2024qualtran} is the most direct structural
precedent for our IR: algorithms compose named \emph{bloqs} with
declared signatures, and the call structure survives down to resource
analysis and surface-code cost modeling.  Its standard library includes
block-encoding constructors, generalized QSP (which subsumes the QSVT
machinery), qubitization, and Hamiltonian-simulation components.
Bloqs can provide symbolic call counts without a complete decomposition,
and Qualtran provides serialization and testing utilities
\cite{qualtran2026interfaces}.
Consequently, neither incomplete decomposition nor serializability alone
distinguishes our design. Our comparison concerns the complete workflow
of restoring an open RIR program, checking replacement interfaces,
capturing resources, and tracking costs before and after partial binding.
pyLIQTR~\cite{obenland2026pyliqtr} generates QSP/QSVT and block-encoding
circuits over Cirq and Qualtran with Clifford+T resource analysis; its
\texttt{clam} layer integrates differential equations \emph{classically},
and no quantum solver families ship with it.
pyqsp~\cite{chuang2021pyqsp} computes QSP phase angles by exact
divide-and-conquer and Newton iteration and is the numerical predecessor
of our phase engine; our engine differs in verifying its computed phases
against the polynomial identity inside the language's validation pass,
rather than acting as an offline numerical tool.
qsp4pde~\cite{kim2026qsp4pde} applies QSP to linear PDEs (advection,
wave, and Poisson equations) in Fourier space with Qiskit-based circuits
and hardware experiments---a single-path solver package, not a language
or an IR.
OpenFermion~\cite{mcclean2020openfermion} plays the analogous input-model
role for quantum chemistry (fermion--qubit mappings, low-rank
decompositions, Trotterization), with no QLSS/QSVT or PDE content.
An external reimplementation study---six published-anchor tasks
reimplemented across these frameworks under a pre-registered
protocol---quantifies the resulting code-volume differences in
Appendix~\ref{app:expressiveness}, keeping language-level expressiveness
and bundled-numerics capability separate (our own degree-limited phase
synthesis is an implementation matter behind a protocol seam, recorded
there as such).
These libraries confirm that fault-tolerant scientific computing demands
structure-preserving composition. The case studies below evaluate our
specific combination of persisted open modules, compatible
realizations, scientific arithmetic, and solver interfaces; these
comparisons do not establish that other frameworks cannot be extended
to express the same computations.

\paragraph{Commercial and industrial platforms.}
A second class of counterparts is commercial.
UnitaryLab~\cite{unitarylab2025platform}, launched in 2025 by the
Shanghai Jiao Tong University team behind the Schr\"odingerization
program~\cite{jin2024schrodingerization,hu2024schrocircuits}, positions
itself as a dedicated quantum scientific computing platform: an
agent-driven web studio over a proprietary simulator SDK, with an
MIT-licensed algorithm library spanning HHL, a QSVT linear-system solver,
variational and adiabatic alternatives, several Hamiltonian-simulation
kernels, and Schr\"odingerization-based heat- and advection-equation
solvers---the closest commercial counterpart to our ODE/PDE coverage.  It
is a product stack rather than a language: programs are
platform-internal, and no open-program artifact, oracle contract, or
resource estimation is surfaced in its public materials.
Classiq~\cite{classiq2026platform} industrializes the synthesis route: the
statically typed Qmod language models algorithms as constraint-annotated
function hierarchies that a closed synthesis engine compiles to
hardware-constrained circuits, and its example library includes QSVT
matrix inversion, HHL, quantum differential equations, and GQSP---but
abstraction ends at synthesis time; there is no open oracle, no late
linking, and no serializable open program.  On the axis we deliberately
reserve, the Microsoft Azure Quantum Resource
Estimator~\cite{microsoft2026azureqe} is the strongest industrial answer:
language-agnostic fault-tolerant estimation from Q\#, QIR, or logical
counts through to physical-qubit and runtime estimates under
error-correction
assumptions~\cite{beverland2022assessing}.  Vertical scientific products
such as Quantinuum's InQuanto~\cite{quantinuum2026inquanto} (quantum
chemistry, with variational and phase-estimation methods and
quantum-error-corrected experiments in its recent results) and algorithm
companies such as
Phasecraft~\cite{phasecraft2026} deliver domains and results rather than
languages, and execution infrastructure such as NVIDIA's
CUDA-Q~\cite{nvidia2026cudaq} and Amazon
Braket~\cite{amazonbraket2026sdk} targets hybrid workloads at the circuit
level. These systems motivate comparison by research task and artifact
boundary, with physical resource estimation and backend performance
treated as separate capabilities.

\paragraph{Research codebases for scientific algorithms.}
Individual algorithm implementations---LCHS, Schr\"odingerization, Carleman
pipelines, QLSS prototypes---exist as research code accompanying
publications~\cite{an2026lchs,jin2024schrodingerization,liu2021pnas,
liu2023carleman,an2026lchs}.  Each typically fixes one method, one input model, and one
simulation regime. Our framework provides shared access-model interfaces
and explicit adapters for connecting such constructions. Applicability
still depends on each algorithm's mathematical assumptions and output
conventions; a common interface does not make every pair composable.

\begin{table}[t]
  \centering
  \caption{Language and framework comparison against the requirements of
  \secref{sec:qsc-requirements}.  \ok~= provided; \half~= partial (see
  text); \noent~= absent.  Claims are referenced to the systems'
  publications and documentation.}
  \label{tab:qpl-comparison}
  \small
  \setlength{\tabcolsep}{4pt}\begin{tabular}{@{}p{2.4cm}cccccc@{}}
    \toprule
    & \textbf{Quipper} & \textbf{Scaffold} & \textbf{Q\#} & \textbf{ProjectQ} & \textbf{Qiskit} & \textbf{\pqlang} \\
    \midrule
    R1 deferred oracles & \noent$^{a}$ & \noent & \half$^{b}$ & \half$^{c}$ & \noent & \ok \\
    R2 realization independence & \noent & \noent & \half & \noent$^{c}$ & \half & \ok \\
    R3 structure-preserving IR & \half & \half & \half & \noent & \noent & \ok \\
    R4 adjoint/control tracked & \half & \noent & \ok & \noent & \noent & \ok \\
    R5 symbolic structure & \half & \half & \half & \noent & \noent & \ok \\
    R6 promise provenance & \noent & \noent & \noent & \noent & \noent & \ok \\
    R7 readout separation & \half & \half & \noent & \noent & \noent & \ok \\
    R8 library extensibility & \half & \half & \half & \half & \half & \ok \\
    \midrule
    Scientific primitives$^{d}$ & \noent & \noent & \half & \noent & \half & \ok \\
    ODE/PDE solver families & \noent & \noent & \noent & \noent & \noent & \ok \\
    Resource estimation & \half$^{e}$ & \ok & \ok & \ok & \half & \half$^{f}$ \\
    \bottomrule
    \multicolumn{7}{@{}p{14.3cm}@{}}{\scriptsize $^a$Boxes abstract over generated subcircuits, but there is no open-program state.  $^b$Operations are passed as parameters and remain abstract to the algorithm, but bodies are concrete and no link-time openness or paradigm contracts exist.  $^c$The extensible engine chain can emulate oracles at execution time, but not in a serializable program artifact.  $^d$Block-encoding algebra,
    QSVT/QSP engine, sparse-access and QRAM input models, reversible
    fixed-point arithmetic with a math-function compiler.
    $^e$Resource analyses exist as external studies over generated
    circuits~\cite{scherer2017resource}; no integrated estimator.
    $^f$Logical gate and QRAM-access estimation is implemented
    (\secref{sec:resources-summary}); physical resource estimation remains
    outside scope (\secref{sec:reserved}).}
  \end{tabular}
\end{table}

\begin{table}[t]
  \centering
  \caption{Algorithm and solver coverage: \pqlang{} against the major
  frameworks' standard libraries.  The scoring counts each framework's
  \emph{own} standard library only; ecosystem-level implementations
  (Qualtran, pyLIQTR, PennyLane's subroutine templates, Classiq's
  example library, UnitaryLab's algorithm library) are surveyed
  separately in Table~\ref{tab:platform-comparison} and the text, so
  that no side is credited or discounted by the choice of boundary.
  \ok~= provided as a composable, contract-checked library component;
  \half~= demonstration-level implementation; \noent~= absent.}
  \label{tab:ode-coverage}
  \small
  \setlength{\tabcolsep}{4pt}\begin{tabular}{@{}p{5.4cm}ccccc@{}}
    \toprule
    \tabhead{Capability} & \textbf{Q\#} & \textbf{ProjectQ} & \textbf{Qiskit} & \textbf{Cirq} & \textbf{\pqlang} \\
    \midrule
    Query/search/amplification & \ok & \ok & \ok & \ok & \ok \\
    Phase/amplitude estimation & \ok & \ok & \ok & \half & \ok \\
    Hamiltonian simulation (Trotter) & \ok & \half & \ok & \half & \ok \\
    Hamiltonian simulation (LCU/Taylor, QSVT) & \noent & \noent & \half & \noent & \ok \\
    Block-encoding algebra & \noent & \noent & \half & \noent & \ok \\
    QSP phase synthesis (self-verifying) & \noent & \noent & \noent & \noent & \ok \\
    Quantum linear systems (QLSS) & \noent & \noent & \half$^{a}$ & \noent & \ok \\
    QRAM/QROM input models (typed resource) & \noent & \noent & \half & \noent & \ok \\
    Reversible fixed-point arithmetic & \ok & \noent & \noent & \noent & \ok \\
    Math-function compiler (Python $\rightarrow$ reversible) & \noent & \noent & \noent & \noent & \ok \\
    ODE: LCHS & \noent & \noent & \noent & \noent & \ok \\
    ODE: Schr\"odingerization & \noent & \noent & \noent & \noent & \ok \\
    ODE: contour decomposition (CBMD) & \noent & \noent & \noent & \noent & \ok \\
    ODE: Carleman (nonlinear) & \noent & \noent & \noent & \noent & \ok \\
    ODE: history systems via QLSS & \noent & \noent & \noent & \noent & \ok \\
    PDE: heat / Fokker--Planck input models & \noent & \noent & \noent & \noent & \ok \\
    PDE: symbolic nonlinear PDEs (QHAM) & \noent & \noent & \noent & \noent & \ok \\
    PDE: Euler equations (QFVM/Roe) & \noent & \noent & \noent & \noent & \ok \\
    \bottomrule
    \multicolumn{6}{@{}p{14.0cm}@{}}{\scriptsize $^a$A demonstration-level HHL circuit shipped historically in Qiskit's (since-retired) algorithms module; it is not part of the current library.}
  \end{tabular}
\end{table}

\begin{table}[t]
  \centering
  \caption{Fault-tolerant algorithm-analysis libraries and commercial
  scientific-computing platforms, compared with \pqlang{} on the axes of
  \secref{sec:qsc-requirements}.  \ok~= provided; \half~= partial (see
  text and footnotes); \noent~= absent.  Claims are referenced to the
  systems' publications, repositories, and documentation; entries for
  closed-source products are limited to what those sources state.}
  \label{tab:platform-comparison}
  \small
  \setlength{\tabcolsep}{3pt}\begin{tabular}{@{}p{3.2cm}cccccc@{}}
    \toprule
    & \textbf{Qualtran} & \textbf{pyLIQTR} & \textbf{PennyLane} & \textbf{Classiq} & \textbf{UnitaryLab} & \textbf{\pqlang} \\
    \midrule
    Open-program oracles (R1/R2) & \half$^{a}$ & \noent & \noent & \half$^{b}$ & \noent & \ok \\
    Structure-preserving IR (R3) & \ok & \noent & \noent & \half & \noent & \ok \\
    Block-encoding primitives & \ok & \ok & \half$^{c}$ & \ok & \half$^{d}$ & \ok \\
    QSP/QSVT phase machinery & \ok & \ok & \half$^{c}$ & \ok & \ok & \ok$^{e}$ \\
    QLSS library components & \noent$^{f}$ & \noent & \noent & \ok & \ok & \ok \\
    ODE/PDE solver families & \noent & \noent$^{g}$ & \noent & \half & \ok & \ok \\
    Fault-tolerant resource estimation & \ok & \ok & \noent & \half & \noent & \half$^{h}$ \\
    Open-source core & \ok & \ok & \ok & \noent$^{i}$ & \half$^{j}$ & \ok \\
    \bottomrule
    \multicolumn{7}{@{}p{14.3cm}@{}}{\scriptsize $^a$Bloqs may specify symbolic call counts without complete decompositions; serialization is supported. The comparison concerns the checked partial-linking workflow, not incompleteness alone (\cite{qualtran2026interfaces}).  $^b$Qmod functions are abstract until synthesis, which concretizes them; no serializable open program.  $^c$Subroutine templates (\texttt{QSVT}, \texttt{GQSP}, \texttt{FABLE}) rather than packaged algorithms; phase angles are supplied externally.  $^d$Block encodings underlie the QSVT linear solver, while the Schr\"odingerization path documents its block-encoding backend as pending.  $^e$Self-verifying: computed phases are checked against the polynomial identity during validation.  $^f$Block-encoding and QSP bloqs exist; no assembled solver.  $^g$The \texttt{clam} layer integrates ODEs classically.  $^h$Logical-level estimator; the physical layer is reserved (\secref{sec:reserved}).  $^i$Closed synthesis engine; the example library is open.  $^j$Proprietary simulator SDK; the algorithm library is MIT-licensed.}
  \end{tabular}
\end{table}

\subsection{Positioning}
\label{sec:positioning}

The comparison is not a claim that prior systems are poorly built; they are
built for a different workload.  Circuit SDKs optimize for NISQ execution
and hardware co-design; typed language lines optimize for semantic safety
of full quantum computation; Q\# optimizes for a vertically integrated
development kit.  \pqlang{} optimizes for the fault-tolerant
scientific-computing workload, where the unit of progress is a
\emph{mathematical algorithm against declared access models}, the unit of
sharing is a \emph{program artifact with open slots and recorded promises},
and the unit of validation is \emph{graduated evidence} rather than a single
simulator's output.  The commercial and analysis-library landscape of
Tables~\ref{tab:qpl-comparison}, \ref{tab:ode-coverage}, and
\ref{tab:platform-comparison} are inventories of documented interfaces in
the surveyed versions, not impossibility results. Missing bundled
components must be distinguished from missing expressiveness.
Our implementation study therefore reports the selected path and any
additional construction required. Qualtran and pyLIQTR remain the closest
comparators for compositional algorithms and logical resource analysis;
physical-cost estimates remain outside our present implementation scope.
